\section{Discussion}
\label{sec:discussion}

\paragraph{The accounting} The question this paper has asked at every level of
the problem is whether a parameter that is not recovered was lost by the
estimator or never present in the data, and the answer has come from one
computation, the Fisher information of the experiment at the truth. The answer
was different at different levels. For the 34 regularly behaved kinetic
parameters of the mechanistic model, the bound said nearly everything is in the
data, and the failures of earlier PINN designs, which had been read as a hard
ceiling of the model, were an estimator gap: a converged classical estimator
reaches the bound, and a Gauss--Newton PINN reaches the classical estimator
where its continuation converges. For APC functionality in the high-WNT
regimes and for the basal rate beside a learned term, the bound said the data
do not contain the parameter, and no estimator change we tried, whether a
different residual, a second-order optimiser, a converged maximum-likelihood
fit or five different constraint parameterisations of the learned term, recovered it. In
both cases the bound also names the remedy: a condition that leaves the \bcat{}
plateau for $\thP$, and a condition that visits the anchor while still
containing the basal term for the learned production edges. Each remedy is
confirmed by fits (for the learned term, fits run before the bound was
computed), against an information-matched control that added the same amount of data without visiting
the relevant region.

\paragraph{Why the bound is the right yardstick for PINNs} A PINN's
parameter-recovery count is a property of the estimator and the data jointly,
and without a yardstick it cannot be attributed to either. The Cram\'er--Rao
bound costs one sensitivity computation (about 720 ODE solves per regime here,
under four minutes on 24 cores). It is local and asymptotic, and both qualifications matter. We therefore
verified that the bound is attained, by Monte Carlo over independent datasets
(Section~\ref{sec:res-mc}), and we reported where it is not exactly attained:
in the two high-WNT regimes the converged estimates scatter less than the bound
for the most weakly informed parameters. The likelihood also has more than one
basin, and the classical estimator from a single start can be trapped (in Early
Adenoma it was); a few starts selected by cost were enough here, but in general
the gap between a fit and the bound can be partly an optimisation gap of the
likelihood itself, not of the PINN. Profile likelihood~\citep{raue2009identifiability,kreutz2013profile,simpson2023pwa}
is the standard non-local alternative; it requires a full re-optimisation of the
other 35 parameters at every point of every profile, and for the questions
asked here (which parameters does this design inform, and which added
condition informs them) a bound computed once and checked once by Monte
Carlo, was sufficient. The practical recommendation is modest. On synthetic benchmarks, compute the
bound of the design before training and report recovery against it: a
parameter whose error exceeds the bound by orders of magnitude is an estimator
problem. On real data the truth is unknown, but the same comparison can be made
between estimators: the bound at the estimate predicts how far two estimators
on the same data should agree, a fit whose parameters differ from a converged
classical fit by many bound-widths has not converged, multistart costs reveal
the multimodality that makes any local bound misleading in a wrong basin, and
for the SGN-PINN the physics residual says whether its continuation reached an
ODE solution.

\paragraph{The estimator side} The Schur-complement Gauss--Newton PINN is not a
new optimiser; it is the observation that freezing a trained PINN's hidden features, so that each
state is linear in its output layer, turns the inverse problem into a separable
nonlinear least-squares problem, for which Gauss--Newton with variable projection is the classical
tool~\citep{golub1973differentiation,golub2003separable,oleary2013varpro}, and
that the block-arrowhead structure of its normal equations makes each step cost
ten small factorisations. Random-feature and extreme-learning-machine
PINNs~\citep{huang2006elm,dong2021local,dong2023inverse} exploit the same
linearity for forward problems; second-order and natural-gradient optimisers
for PINNs~\citep{muller2023ngd,jnini2024gauss} are a related line. What the
formulation adds for inverse problems is that the Schur complement it already
computes is the head-profiled Gauss--Newton curvature, which approximates the
profiled Fisher information when the physics term is satisfied, so the
estimator reports an estimate of its own CRLB at no extra cost. That estimate of uncertainty is only as good as the point
it is computed at: where the continuation stalled, the SGN-PINN reported
confident standard deviations around wrong estimates, which is the general
limitation of any local uncertainty estimate~\citep{transtrum2015perspective}.

\paragraph{The information side, and its relation to prior work} The
degeneracy behind the learned-term result is known: Loman and Baker construct
it for a linear-rate term and note its generality~\citep{loman2025identifiability}.
They also find that monotone, bounded constraints bring parametric
identifiability on par with the fully known model, which is not what our
monotone exact-anchor or oracle-bounded forms achieve under the base design. We
think the difference is the variable they do not treat, the data support of
the learned function: the same forms recover the basal rate the moment one
condition reaches the anchor. The survey of~\citet{philipps2025ude} reports a
learned term recovered with a spurious additive constant and attributes the
mismatch to non-identifiabilities introduced by the network; the compensation
identity, Eq~\eqref{eq:compensation}, says precisely which constant, and the
hybrid Fisher information says which experiment would have removed it. To our
reading the only optimal-experimental-design work for universal differential
equations~\citep{plate2024oed} designs on the joint Fisher information of
network weights and mechanistic parameters without shape constraints, so it has
no notion of designing a condition so that a constraint binds; classical
design for dynamic models~\citep{kreutz2009design,bandara2009oed,apgar2010sloppy,transtrum2012oed,liepe2013maximizing}
optimises a scalar of the Fisher matrix over inputs and sampling times. The
anchor ratio is a cruder, cheaper and constraint-aware surrogate, and the
hybrid Fisher information is its quantitative version. In the language of
adaptive control, both information gaps are failures of persistence of
excitation~\citep{sastry1989adaptive,wang2006learning}: a learned map is
identified only along the trajectory the system follows, and so, we add, is any
constraint imposed on it.

\paragraph{Biological reading} The model is a reduced caricature of the crypt,
but two of its conclusions translate into experimental advice that does not
depend on its details. APC functionality, the quantity that distinguishes the
disease stages, is the parameter that a perturbation design built around
retinoid treatment and WNT/MYC pulses informs least exactly where it matters
most, in the high-WNT states, because \bcat{} saturates the readouts that would
report it; a graded WNT knockdown restores it, while two added conditions that
perturb without lowering \bcat{} do not. And when a regulatory law is learned by a flexible network
rather than assumed, the basal transcription rate of the target gene, a
quantity of direct biological interest, is identified either by a condition that
silences the regulator, which in a positive-feedback network may require
knocking down the feedback as well as the input (the single WNT knockout leaves
\bcat{} at about 1\% of its range because HOXA13 keeps producing it), or by a known
partial knockdown of the basal term itself.

\paragraph{Limitations} Every result is on synthetic data generated by the model
itself, with all seven states observed at 150 time points with noise of 0.1--2.6\% of their mean levels, so
every recovery count is an upper bound on what real measurements of a subset of
states would give, and the regimes are not calibrated to patient data. The
information budget depends strongly on that design (S1~Fig): with only
\bcat{}, the two HOX genes and MYC observed, an efficient estimator would
recover 16--30 of the 34 regular parameters per regime, with \bcat{} and MYC
alone 4--8.5; with ten-fold larger noise and all states observed, 17--32. The accounting itself does
not depend on the design, and it is exactly what would tell a modeller which of
these regimes they are in before any network is trained. The
Cram\'er--Rao analysis is local and asymptotic: in the two high-WNT regimes it
is conservative for the weakest parameters, and the likelihood has more than one
basin. The hybrid experiments other than the full-pipeline
confirmations use an equation-local screen with exact states; the full pipeline
was rerun for two edges and confirms the direction of the effects, with larger
absolute errors than the screen's because state estimation is in the loop. Restart counts move individual hybrid cells by up to fifty points
away from the anchor, so numbers are compared within experiments; the headline
contrast at the anchor is replicated over three disjoint seed sets. As registered, two of the six dose-response
predictions passed (three after rerunning one regime with more restarts) and two
of the five prospective-test predictions (S1~Text), including the registered
falsifier failing in one regime; the second clause of the design rule was
formulated after the failure that revealed it, and the hybrid Fisher information reproduces it but
was computed afterwards. The prescribed protocols for the MYC and HOXA13 edges
require silencing a constitutive promoter, which is materially harder than the
medium change that fixes the RA edges, and the Fisher analysis points to a
cheaper alternative it did not test by fit: a known partial knockdown of a basal
term identifies it without visiting the anchor. Every knockdown here, and that
alternative in particular, assumes the knockdown efficiency is known exactly;
real siRNA efficiencies are uncertain, and an unknown efficiency adds one
parameter per knockdown condition that a design would have to identify too. Finally, no symbolic
structural-identifiability analysis~\citep{villaverde2016structural,chis2011structural}
was run; apart from the $\dP(1-\thP)$ product, every identifiability statement
here is practical.

\paragraph{Conclusion} A recovery count does not say whether to change the
estimator or the experiment; the Cram\'er--Rao bound of the design does, at the
cost of one sensitivity computation. In a seven-state model of colorectal-cancer signalling it
showed that the long-standing recovery ceiling of physics-informed estimators
was numerical, removable by a converged classical fit and by the SGN-PINN where
its continuation converges, that the parameter
defining the disease stage is uninformed by the conventional design in exactly
the states that matter, and that the information a hybrid model's constraint
carries is set by where the experiment takes the regulator. For neural and
hybrid inverse problems in systems biology, the first diagnostic should be the
information budget of the design.
